% GENERATED FILE. Edit canonical lessons, then run npm run generate:books.
% canonical-source-hash: fnv1a64:c6bb74ae9392cac2
% canonical-lessons: LA-C267-pinguis, LA-C267-callidus, LA-C267-praecipuus, LA-C267-pristinus, LA-C267-proximus

\chapter{Fat, Clever, Special, Former, Nearest}
\label{ch:la-fat-clever-special-former-nearest}
\hlchaptermodality{267}

\begin{chapteropening}
\textbf{By the end of this chapter:} \emph{I can say fat, clever, special, former and nearest.}

Complete the last lesson of chapter 267: I can say fat, clever, special, former and nearest.
\end{chapteropening}

\section[pinguis, pingue]{pinguis, pingue — fat}

\label{lesson:LA-C267-pinguis}



\begin{quote}
Before the new one: say the Latin for fair, then the Latin for unjust.
\end{quote}

\subsection*{You'll want to know: pinguis, pingue}
\textbf{pinguis, pingue} — \textquotedblleft{}fat\textquotedblright{}.

\begin{cousinweb}
\textbf{pinguis} has no sure origin; it is often linked with a Proto-Indo-European root \emph{*peyH-}, \textquotedblleft{}fat, to swell\textquotedblright{} (EtymDB 2.1), or with Greek \emph{pakhys}, \textquotedblleft{}thick\textquotedblright{} (Lewis \& Short), but neither link is settled.
\end{cousinweb}

\begin{grammarlens}[title={What you've built}]
One more word for this chapter.
\end{grammarlens}

\subsection*{Your turn}
\begin{itemize}
\raggedright
  \item \emph{Say it:} \emph{pinguis, pingue}
  \item \emph{Say it:} \emph{pinguis, pingue}, once more
  \item \emph{Say it:} say \emph{pinguis, pingue}
  \item \emph{Recall:} say the Latin for fair, then the Latin for unjust, then say \emph{pinguis, pingue} again
\end{itemize}

\begin{tcolorbox}[breakable,colback=teal!4,colframe=teal!35!black,title={Before you move on}]
What does pinguis, pingue mean? (\textquotedblleft{}Fat\textquotedblright{}.) Say it once more.
\end{tcolorbox}
\section[callidus, callida, callidum]{callidus, callida, callidum — clever, cunning}

\label{lesson:LA-C267-callidus}



\begin{quote}
Before the new one: say the Latin for unjust, then the Latin for fat.
\end{quote}

\subsection*{You'll want to know: callidus, callida, callidum}
\textbf{callidus, callida, callidum} — \textquotedblleft{}clever, cunning\textquotedblright{}.

\begin{cousinweb}
\textbf{callidus} is built on \emph{calleō}, \textquotedblleft{}to be thick-skinned, to be practised\textquotedblright{} (Lewis \& Short), from \emph{callum}, \textquotedblleft{}hard skin\textquotedblright{}: hardened by experience, so shrewd. English \emph{callous} comes from the same family.
\end{cousinweb}

\begin{grammarlens}[title={What you've built}]
One more word for this chapter.
\end{grammarlens}

\subsection*{Your turn}
\begin{itemize}
\raggedright
  \item \emph{Say it:} \emph{callidus, callida, callidum}
  \item \emph{Say it:} \emph{callidus, callida, callidum}, once more
  \item \emph{Say it:} say \emph{callidus, callida, callidum}
  \item \emph{Recall:} say the Latin for unjust, then the Latin for fat, then say \emph{callidus, callida, callidum} again
\end{itemize}

\begin{tcolorbox}[breakable,colback=teal!4,colframe=teal!35!black,title={Before you move on}]
What does callidus, callida, callidum mean? (\textquotedblleft{}Clever, cunning\textquotedblright{}.) Say it once more.
\end{tcolorbox}
\section[praecipuus, praecipua, praecipuum]{praecipuus, praecipua, praecipuum — special, chief}

\label{lesson:LA-C267-praecipuus}



\begin{quote}
Before the new one: say the Latin for fat, then the Latin for clever.
\end{quote}

\subsection*{You'll want to know: praecipuus, praecipua, praecipuum}
\textbf{praecipuus, praecipua, praecipuum} — \textquotedblleft{}special, chief\textquotedblright{}.

\begin{cousinweb}
\textbf{praecipuus} is built on \emph{praecipiō}, \textquotedblleft{}to take beforehand\textquotedblright{} (Lewis \& Short), from \emph{prae}, \textquotedblleft{}before\textquotedblright{}, and \emph{capiō}, \textquotedblleft{}to take\textquotedblright{}: taken first, so chief.
\end{cousinweb}

\begin{grammarlens}[title={What you've built}]
One more word for this chapter.
\end{grammarlens}

\subsection*{Your turn}
\begin{itemize}
\raggedright
  \item \emph{Say it:} \emph{praecipuus, praecipua, praecipuum}
  \item \emph{Say it:} \emph{praecipuus, praecipua, praecipuum}, once more
  \item \emph{Say it:} say \emph{praecipuus, praecipua, praecipuum}
  \item \emph{Recall:} say the Latin for fat, then the Latin for clever, then say \emph{praecipuus, praecipua, praecipuum} again
\end{itemize}

\begin{tcolorbox}[breakable,colback=teal!4,colframe=teal!35!black,title={Before you move on}]
What does praecipuus, praecipua, praecipuum mean? (\textquotedblleft{}Special, chief\textquotedblright{}.) Say it once more.
\end{tcolorbox}
\section[prīstinus, prīstina, prīstinum]{prīstinus, prīstina, prīstinum — former, of old}

\label{lesson:LA-C267-pristinus}



\begin{quote}
Before the new one: say the Latin for clever, then the Latin for special.
\end{quote}

\subsection*{You'll want to know: prīstinus, prīstina, prīstinum}
\textbf{prīstinus, prīstina, prīstinum} — \textquotedblleft{}former, of old\textquotedblright{}.

\begin{cousinweb}
\textbf{prīstinus} belongs with \emph{prīscus}, \textquotedblleft{}ancient\textquotedblright{} (Lewis \& Short), and with \emph{prior}, \textquotedblleft{}earlier\textquotedblright{}, all built on what is usually reconstructed as \emph{*pri-}, \textquotedblleft{}before\textquotedblright{}.
\end{cousinweb}

\begin{grammarlens}[title={What you've built}]
One more word for this chapter.
\end{grammarlens}

\subsection*{Your turn}
\begin{itemize}
\raggedright
  \item \emph{Say it:} \emph{prīstinus, prīstina, prīstinum}
  \item \emph{Say it:} \emph{prīstinus, prīstina, prīstinum}, once more
  \item \emph{Say it:} say \emph{prīstinus, prīstina, prīstinum}
  \item \emph{Recall:} say the Latin for clever, then the Latin for special, then say \emph{prīstinus, prīstina, prīstinum} again
\end{itemize}

\begin{tcolorbox}[breakable,colback=teal!4,colframe=teal!35!black,title={Before you move on}]
What does prīstinus, prīstina, prīstinum mean? (\textquotedblleft{}Former, of old\textquotedblright{}.) Say it once more.
\end{tcolorbox}
\section[proximus, proxima, proximum]{proximus, proxima, proximum — nearest, next}

\label{lesson:LA-C267-proximus}



\begin{quote}
Before the new one: say the Latin for special, then the Latin for former.
\end{quote}

\subsection*{You'll want to know: proximus, proxima, proximum}
\textbf{proximus, proxima, proximum} — \textquotedblleft{}nearest, next\textquotedblright{}.

\begin{cousinweb}
\textbf{proximus} is usually traced to Proto-Indo-European \emph{*prokwismmos}, \textquotedblleft{}nearest\textquotedblright{} (EtymDB 2.1), the superlative belonging to \emph{prope}, \textquotedblleft{}near\textquotedblright{}; Lewis \& Short list it under \emph{propior}, \textquotedblleft{}nearer\textquotedblright{}.
\end{cousinweb}

\begin{grammarlens}[title={What you've built}]
One more word for this chapter.
\end{grammarlens}

\subsection*{Your turn}
\begin{itemize}
\raggedright
  \item \emph{Say it:} \emph{proximus, proxima, proximum}
  \item \emph{Say it:} \emph{proximus, proxima, proximum}, once more
  \item \emph{Say it:} say \emph{proximus, proxima, proximum}
  \item \emph{Recall:} say the Latin for special, then the Latin for former, then say \emph{proximus, proxima, proximum} again
\end{itemize}

\begin{tcolorbox}[breakable,colback=teal!4,colframe=teal!35!black,title={Before you move on}]
What does proximus, proxima, proximum mean? (\textquotedblleft{}Nearest, next\textquotedblright{}.) Say it once more.
\end{tcolorbox}
